% \overfullrule=1mm

%%% The following setting enable figures not to float too much

\newcommand*\hw[1]{\small\fontfamily{augie}\selectfont {#1}}
\newcommand*\circled[1]{
     \tikz[baseline=(char.base)]{
          \node[shape=circle,draw,inner sep=0.5pt,fill=black] (char) {{\color{white} #1}};
     }\!\!
}

\setcounter{topnumber}{2}
\setcounter{bottomnumber}{2}
\setcounter{totalnumber}{4}
\renewcommand{\topfraction}{0.85}
\renewcommand{\bottomfraction}{0.85}
\renewcommand{\textfraction}{0.15}
\renewcommand{\floatpagefraction}{0.7}
%%%%%%%%%%%%%%%%%%%

%%% Macros useful almost everywhere
\newcommand{\fnrestriction}[1]{\lvert_{#1}}
\usepackage{xargs}
\usepackage[normalem]{ulem} % underline command breaks over line ends
% \usepackage[UKenglish]{babel}
% \usepackage{bm}  some strange fonts needed somewhere

\usepackage{xifthen}        % for conditional commands
\newcommand{\ifempty}[3]{%
  \ifthenelse{\isempty{#1}}{#2}{#3}%
}

\DeclareGraphicsExtensions{%
  .png,.PNG,%
  .pdf,.PDF,%
  .jpg,.mps,.jpeg,.jbig2,.jb2,.JPG,.JPEG,.JBIG2,.JB2
}

\newcommand{\singleton}[1]{\textbf{#1}}

%%% My flags
\newif\ifemi
%%% My flags end

\newcommandx{\preprint}[3][1=preprint,2=Springer]{
  \ifempty{#1}{}{
	 \ \\[1em]\noindent
	 \textbf{Disclaimer}
    The published version of this paper is~\cite{#1} (\copyright\ #2).
  }
}

%% Meta-comments
\usepackage{fixme}
\fxusetheme{color}
\FXRegisterAuthor{eM}{aeM}{\color{orange} {\underline{eM}}}
%\usepackage[disable]{todonotes}
\newcommand{\eMtodo}[1]{\todo[fancyline,color=orange!10]{\tiny eM: \color{Navy}#1}}
\usepackage[most]{tcolorbox}
\newtcolorbox{markbox}{
  enhanced,
  breakable,
  size=minimal,
  parbox=false,
  after={\par},
  before upper={\indent},
  colback=white,
  overlay = {
	 \draw[line width=2pt]
	 (frame.north east) -| ([xshift=3mm]frame.east) |-(frame.south east);
  },
  overlay first={\draw[line width=2pt] (frame.north east) -| ([xshift=3mm]frame.south east);},
  overlay middle={\draw[line width=2pt] ([xshift=3mm]frame.north east) -- ([xshift=3mm]frame.south east);},
  overlay last={\draw[line width=2pt] ([xshift=3mm]frame.north east) |- (frame.south east);},
}

\usepackage{todonotes}
\newcommand{\eMcomm}[2][check]{%
  \ifthenelse{\equal{#1}{new}}{{\color{red}#2}}{%
	 \ifthenelse{\equal{#1}{changed}}{{\color{teal}{#2}}}{%
		\ifthenelse{\equal{#1}{rm}}{\todo[color=black!3]{\tiny eM: removed\\'{#2}'}}{%
		  \todo[color=orange!20]{\tiny eM: \color{NavyBlue}#1}%
		  {\color{OliveGreen}{#2}}%
		}%
	 }%
  }%
}
\newcommand{\CLPcomm}[2][check]{%
  \ifthenelse{\equal{#1}{new}}{{\color{red}#2}}{%
	 \ifthenelse{\equal{#1}{changed}}{{\color{teal}{#2}}}{%
		\ifthenelse{\equal{#1}{rm}}{\todo[color=black!3]{\tiny CLP: removed\\'{#2}'}}{%
		  \todo[color=orange!20]{\tiny CLP: \color{NavyBlue}#1}%
		  {\color{OliveGreen}{#2}}%
		}%
	 }%
  }%
}
\newenvironment{eMenv}[1][new]{
  \ifthenelse{\equal{#1}{new}}{\color{red}}{%
	 \ifthenelse{\equal{#1}{changed}}{\color{teal}}{%
		\todo[color=orange!20]{\tiny eM: \color{NavyBlue}#1}%
		\color{OliveGreen}%
	 }%
   }%
  \vspace{2em}\hrule\vspace{.5em}
}{\vspace{.5em}\hrule\vspace{2em}}
\newcommand{\revision}[3]{% example: \revision{R3 asks to change an example with x}{The new text for x}{text of the old example}
  {\textcolor{NavyBlue}{#2}}\xspace\todo[fancyline,color=orange!10]
}
\newcommand{\nota}[2][TODO]{::\textcolor{red}{#1}:: \color{cyan} #2}
\newcommand{\okey}{\dSmiley}
\newcommand{\yeko}{\dXey}
\newcommand{\hidden}[1]{}
\newcommand{\hide}[1]{}

\newcommand{\cf}[2]{
  \fontsize{#1}{#1}{\selectfont{#2}}
}

\makeatletter
\newcommand{\dolist}[2]{%
  %%% \dolist{+}{a,b,c} = a+b+c
  \def\nextitem{\def\nextitem{#1}}%
  \@for \el:=#2\do{\nextitem\textbf{\el}}%
}

\newcommand{\domathlist}[2]{%
  %%% like dolist, in math mode
  \def\nextitem{\def\nextitem{\ensuremath{#1}}}%
  \@for \el:=#2\do{\ensuremath{\nextitem}\textbf{\el}}%
}

\def\mktest#1{
  \def\transform##1+##2+##3{##1 piu' ##2 * ##3\penalty0}%
  \do{\expandafter\transform#1}%
}
\def\mksubscript#1{
  \def\transform##1[##2]{##1_{##2}\penalty0}%
  \do{\expandafter\transform#1}%
}
\newcommand{\mapcmd}[3][{, }]{%
  %%% \mapcmd[\&]{\emph}{uno,due,tre}
  \def\nextitem{\def\nextitem{#1}}%
  \@for \el:=#3\do{\nextitem{#2{\el}}}%
}
\makeatother


\ifemi
\usepackage{showlabels}
\renewcommand{\showlabelsetlabel}[1]{
  {\textcolor{ForestGreen}{
		\showlabelfont{\tiny #1}}
  }
}
\newcommand{\emi}[2]{
  \marginpar{\fcolorbox{red}{shadecolor}{\cf{#1}{{#2}}}}
}
\newcommand{\emic}[2]{\par
  \fcolorbox{red}{shadecolor}{\parbox{\linewidth}{ 
      \color{gray}
      \begin{description}
      \item[{\color{blue} #2}]{\sf #1}
      \end{description}}}
}
\else
\newcommand{\emi}[2]{}
\newcommand{\emic}[2]{}{}
\fi

\newcommand{\marginnote}[1]{
  {\makebox[0pt]{\color{orange}}}
  \marginpar{\parbox{2cm}{\flushleft \tiny \color{orange}{#1}}}
}
\newcommand{\side}[2][6]{
  \marginpar{\cf{#1}{\hl{#2}}}
}
\newcommand{\HSLtag}{\scalebox{1.25}{%
  \begin{tikzpicture}
  \draw (0.1,0.2) -- (0.2,0.115) -- (0.3,0.2) ;
  %
  \draw (0.1,0.2) -- (0.15,0.1) ;
  \draw (0.3,0.2) -- (0.25,0.1) ;
  %
  \draw (0.12,0.1) -- (0.2,0.05) -- (0.28,0.1) ;
  %
  \draw (0.2,0) circle (0.12) ; 
  \draw (0.16,0.04) circle (0.01) ;
  \draw (0.24,0.04) circle (0.01) ;
  \draw (0.15,-0.07) -- (0.25,-0.07) ;
  \end{tikzpicture}
}}
\newcommand{\HSLbox}{\rule{1ex}{1ex}}
\newcommand{\hsl}[1][]{\par
  {\color{red}\vbox{\medskip\noindent\hrulefill \\[5pt]
  \HSLtag \hspace{\stretch{1}}\ifempty{#1}{HIC SUNT
  LEONES}{{#1}\hspace{\stretch{1}}} \HSLtag \\ \smallskip\noindent\hrulefill \\}}\par
}
%%% Meta-comments end


%%% Thanking projects

\newcommand{\tnxbehapi}[1][partly]{%
  Research {#1} supported by the EU H2020 RISE programme under the
  Marie Sk{\l}odowska-Curie grant agreement No 778233%
}

\newcommand{\tnxastra}[1][partly]{%
  Research {#1} supported by the PRO3 MUR project
  Software Quality, and PNRR MUR project VITALITY (ECS00000041), Spoke
  2 ASTRA - Advanced Space Technologies and Research Alliance%
}

\newcommand{\tnxindam}[1][partly]{%
  Research {#1} supported by INdAM as members of GNCS (Gruppo
  Nazionale per il Calcolo Scientifico)%
}

\newcommand{\tnxpro}[1][Work partially funded]{%
   #1 by the PRO3 MUR project Software Quality%
}

\newcommand{\tnxgssi}{%
  The authors acknowledge the support of the PRO3 MUR project Software Quality,
  and PNRR MUR project VITALITY (ECS00000041), Spoke 2 ASTRA - Advanced Space Technologies and Research Alliance%
}
%%% Thanking projects end


%%% Maths & logic
\newcommand{\proofcase}[1]{\ \newline\noindent\fbox{#1}}
\newcommand{\Set}[1]{\left\{\,#1\,\right\}}
\newcommand{\SET}[1]{\big\{\,#1\,\big\}}
\def\colorFun{\color{orange}}
\newcommand{\noarg}{}
\newcommand{\mkfun}[4][\colorFun]{
  \newcommand{#2}[1][#4]{%
    {#1\ensuremath{\mathsf{#3}}}%
    \ifempty{##1}{\noarg}{%
      ({##1})}%
  }%
}

\mkfun{\head}{hd}{}
\mkfun{\tail}{tl}{}

\newcommand{\mkuop}[4][\colorFun]{
  \newcommand{#2}[1][#4]{
    {#1\textsf{#3}}
    \ifempty{##1}{}{
      \, {##1}}
  }
}
\newcommand{\mkset}[1]{\big\{ {#1} \big\}} %%% formerly called \ASET
\newcommand{\partsof}[1]{2^{#1}} %%% formerly called \ASET
\newcommand{\setof}[1]{\mkset{#1}}
\newcommand{\setdef}[2]{\mkset{#1 \sst #2}}
\newcommand{\card}[1]{\|{#1}\|}
\newcommand{\sst}{\;\big|\;}
\newcommand{\qst}{\;\colon\;} %such that
\newcommand{\dom}[1]{\operatorname{dom} {#1}}
\newcommand{\cod}[1]{\operatorname{cod} {#1}}
\newcommand\myenddef{~\hfill \ensuremath{\diamond}}
\newcommand{\conf}[1]{\ensuremath{\langle {#1} \rangle}}
\newcommand{\tuple}[1]{\conf{#1}}
\newcommand{\acsconf}[2]{\conf{{#1} \ ; \ {#2}}}
%\renewcommand{\vec}[1]{\overset{\to}{#1}}
\newcommand{\ith}[2]{\vec{#1}[{#2}]}
\newcommand{\emptyword}{\varepsilon}
\newcommand{\aword}{\varphi}
\newcommand{\langof}[1]{\mathcal{L}(#1)}
\newcommand{\sem}[2][]{\mbox{\ensuremath{\llbracket{#2}\rrbracket_{#1}}}}
\newcommand{\qmmdef}{\quad \mmdef \quad}
\newcommand{\qqmmdef}{\qquad \mmdef \qquad}
\newcommand{\qqand}[1][and]{\qquad\text{#1}\qquad}
\newcommand{\qand}[1][and]{\quad\text{#1}\quad}
\newcommand{\avalue}[1][v]{
  \ensuremath{
	 \mathtt{#1}
  }
}
\newcommand{\verum}{\avalue[true]}
\newcommand{\falsus}{\avalue[false]}
\newcommand{\nat}{\mathbb{N}}
\newcommand{\upd}[3]{{#1}[{#2} \mapsto {#3}]}
\newcommand{\grmeq}{\; ::= \;}
\newcommand{\bnfdef}{\ ::=\ }
\newcommand{\grmor}{\; \big| \;}
\newcommand{\bnfmid}{\;\ \big|\ \;}
\newcommand{\sep}{\;\bnfmid\;}
\newcommand{\subs}[2]{\left[^{#1}/_{#2}\right]}
%%% Maths & logic end



%%% Typographic style
\newcommand{\defrule}[1]{%
  % \hypertarget{rule:{#1}}
  {%
    \text{\scriptsize[\textbf{#1}]}%
  }%
}
\newcommand{\ie}{\text{i.e.,}\xspace}
\newcommand{\cfw}{\text{cf.}\xspace}
\newcommand{\eg}{\text{e.g.,}\xspace}
\newcommand{\aka}{\text{a.k.a.}\xspace}
\newcommand{\mypar}[1]{\noindent\paragraph{\textbf{#1}}\ }
\newcommand{\secref}[1]{\S~\ref{#1}}
\newcommand{\figref}[1]{Fig.~\ref{#1}}
\newcommand{\exref}[1]{Ex.~\ref{#1}}
\newcommand{\defref}[1]{Definition~\ref{#1}}
\newcommand{\exeref}[1]{Exercise~\ref{#1}}
\newcommand{\rmkref}[1]{Remark~\ref{#1}}
\newtheorem{rmk}{Remark}{\bfseries}{\rmfamily}
\newtheorem{exe}{Exercise}{\bfseries}{\rmfamily}
\newcommand{\sidebyside}[2]{
  \begin{tabular}{ll}
    \begin{minipage}{.5\linewidth} {#1}  \end{minipage}
    &
    \begin{minipage}{.5\linewidth} {#2}  \end{minipage}
  \end{tabular}
}
\newcommand{\eqdef}{\ \triangleq\ }
\newcommand{\mmdef}{\eqdef}
%%% Typographic style end



%%% Arrows
%%%%%
%% Extensible dashed arrow
%%%%%
\makeatletter
\newcommand*{\da@rightarrow}{\mathchar"0\hexnumber@\symAMSa 4B }
\newcommand*{\da@leftarrow}{\mathchar"0\hexnumber@\symAMSa 4C }
\newcommand*{\xdashrightarrow}[2][]{%
  \mathrel{%
    \mathpalette{\da@xarrow{#1}{#2}{}\da@rightarrow{\,}{}}{}%
  }%
}
\newcommand{\xdashleftarrow}[2][]{%
  \mathrel{%
    \mathpalette{\da@xarrow{#1}{#2}\da@leftarrow{}{}{\,}}{}%
  }%
}
\newcommand*{\da@xarrow}[7]{%
  % #1: below
  % #2: above
  % #3: arrow left
  % #4: arrow right
  % #5: space left 
  % #6: space right
  % #7: math style 
  \sbox0{$\ifx#7\scriptstyle\scriptscriptstyle\else\scriptstyle\fi#5#1#6\m@th$}%
  \sbox2{$\ifx#7\scriptstyle\scriptscriptstyle\else\scriptstyle\fi#5#2#6\m@th$}%
  \sbox4{$#7\dabar@\m@th$}%
  \dimen@=\wd0 %
  \ifdim\wd2 >\dimen@
    \dimen@=\wd2 %   
  \fi
  \count@=2 %
  \def\da@bars{\dabar@\dabar@}%
  \@whiledim\count@\wd4<\dimen@\do{%
    \advance\count@\@ne
    \expandafter\def\expandafter\da@bars\expandafter{%
      \da@bars
      \dabar@ 
    }%
  }%  
  \mathrel{#3}%
  \mathrel{%   
    \mathop{\da@bars}\limits
    \ifx\\#1\\%
    \else
      _{\copy0}%
    \fi
    \ifx\\#2\\%
    \else
      ^{\copy2}%
    \fi
  }%   
  \mathrel{#4}%
}
\makeatother

\def\ured{\rightarrow}
\def\tred#1{\overstackrel{\rightarrowfill}{#1}}
\def\iured#1{\rightarrow_{\descr{\scriptstyle #1}}}
\def\itred#1#2{\mathrel{{\overstackrel{\rightarrowfill}{#2}}_{\descr
      {\scriptstyle #1}}}}
\def\sytr#1{\overstackrel{\longmapsto}{#1}}
\def\tr#1{\overstackrel{\rightarrowfill}{#1}}
\def\actr#1#2{\overstackrel{\rightarrowfill}{\act {#1}{#2}}}
\def\iactr#1#2#3{\overstackrel{\rightarrowfill}{\act {#2}{#3}}_{\descr
    {\scriptstyle #1}}}
\def\dlarrow#1#2{\overunderstackrel{\rightarrowfill}{#1}{#2}}
\def\ddlarrow#1#2{\overunderstackrel{\longmapsto}{#1}{#2}}
\def\Dlarrow#1#2{\overunderstackrel{\Longrightarrow}{#1}{#2}}
\def\dmtarrow#1#2{\overunderstackrel{\longmapsto}{#1}{#2}}
\def\mtarrow#1{\overstackrel{\longmapsto}{#1}}
\newcommand{\mylarrow}[1]{\overstackrel{\rightarrowfill}{#1}}

\newcommand{\partarrow}[2]{\overunderstackrel{\rightarrowfill}{\ \ \mbox{\scriptstyle $#1$}\ \ }{\ \mbox{\scriptstyle $#2$}\ }\!\!\!>}
\newcommand{\arw}[1]{\overstackrel{\rightarrowfill}{#1}}
\newcommand{\newpartarrow}[2]{\partarrow{#1}{}}
\newcommand{\spartarrow}[1]{\overstackrel{\rightarrowfill}{\ \ \mbox{\scriptstyle $#1$}\ \ \ }\!\!\!\!\!>}
\newcommand{\trlts}[1]{\overstackrel{\rule[.8mm]{7mm}{.2mm}}{#1}\!\!\blacktriangleright}
\newcommand{\trred}{\overstackrel{\rule[.8mm]{7mm}{.2mm}}\!\!\vartriangleright}
\newcommand{\arro}[2][{}]{\xrightarrow[{#1}]{\ {#2}\ }}
%%% Arrows end

%%% Others
\newcommand{\squo}[1]{\lq {#1}\rq}
\newcommand{\quo}[1]{\lq\lq {#1}\rq\rq}
\def\finex{{\unskip\nobreak\hfil
\penalty50\hskip1em\null\nobreak\hfil$\diamond$
\parfillskip=0pt\finalhyphendemerits=0\endgraf}}
\newenvironment{myex}{\begin{example}\it}{\finex\end{example}}
\newcommand{\asort}[1][s]{\mathtt{#1}}
\newcommand{\typeof}[2]{\ensuremath{
	 \vdash \varid[{#1}] \colon \asort[{#2}]
  }\xspace
}
%%% Others end


%%% Some colors
\definecolor{shadecolor}{rgb}{1,0.99,0.9}
\definecolor{bg}{rgb}{0.95,0.95,0.95}
\newcommand{\newgreen}{green!50!blue!100}
%%% Some colors end


%%%%%%%%%%%%%%%%%%%%%%%%%%%%%%%%%
%%%%           ABC           %%%%
%%%%%%%%%%%%%%%%%%%%%%%%%%%%%%%%%
\newcommand{\abc}{\textsf{AbC}\xspace}
\newcommand{\abcST}{\textsf{ABeT}\xspace}
\newcommand{\abcattr}[1][a]{\textsf{#1}}
\newcommand{\abccond}[1][\rho]{#1}
\newcommand{\abcval}[1][v]{\texttt{#1}}
\newcommand{\abcvar}[1][x]{#1}
\def\colorExp{\color{NavyBlue}}
\def\colorTuple{\color{Black}}
\newcommand{\abcexp}[1][e]{\colorExp #1}
\newcommandx{\abctuple}[1][1 = t]{\llparenthesis{#1}\rrparenthesis}
\newcommandx{\abcget}[2][1=a,2={id},usedefault=@]{
  \ptp[{#1}]{\colorOp .}\abcattr[{#2}]
}
\newcommandx{\abcptp}[2][1=a,2=\abccond,usedefault=@]{\ifempty{#1}{}{\ptp[{#1}] \ifempty{#2}{}{{\colorOp \shortmid}}} {#2}}
\newcommandx{\abcint}[6][1=a,2=\abccond,3=e,4=e',5=b,6=\abccond',usedefault=@]{
  \abcptp[{#1}][{#2}]
  \ {\colorOp \xrightarrow{\scriptstyle \abcexp[#3]\quad\abcexp[#4]}}\ 
  \abcptp[{#5}][{#6}]
}
\newcommandx{\mkabcint}[8][3=a,4=\abccond,5=e,6=e',7=b,8=\abccond',usedefault=@]{
  \node[bblock, #1] (#2) {$\abcint[{#3}][{#4}][{#5}][{#6}][{#7}][{#8}]$};
}

\tikzset{
    abccallout/.style={
      fill=green!10,
		opacity=.5,
		overlay,
		align=center,
      cloud callout,
		cloud puffs=15,
		aspect=2.5,
		cloud ignores aspect,
		cloud puff arc=100,
		shading=ball
    }
  }

\newcommandx{\abcP}[6][1=P,2=K,3=.1cm,4=1cm,5=north east,6=proc,usedefault=@]{
  \begin{tikzpicture}
	 \node[fill=blue!10, shape=circle] (#6) {$\p[#1]$};
	 \node[abccallout, above = #3 of #6, xshift=#4, callout absolute pointer={(#6.#5)}] {$#2$}
	 ;	 
	 \draw[decorate,decoration={expanding waves,angle=7,segment length = .05cm}] (#6.east) -- ++(.5cm,0)
	 ;
  \end{tikzpicture}
}


\ExplSyntaxOn
\NewDocumentCommand{\ucgreek}{m}
 {
  \str_case:nn { #1 }
   {
    {A}{\mathrm{A}}
    {B}{\mathrm{B}}
    {C}{\Sigma}
    {D}{\Delta}
    {E}{\mathrm{E}}
    {F}{\Phi}
    {G}{\Gamma}
    {H}{\mathrm{H}}
    {I}{\mathrm{I}}
    {J}{\Theta}
    {K}{\mathrm{K}}
    {L}{\Lambda}
    {M}{\mathrm{M}}
    {N}{\mathrm{N}}
    {O}{\mathrm{O}}
    {P}{\Pi}
    {Q}{\mathrm{X}}
    {R}{\mathrm{P}}
    {S}{\Sigma}
    {T}{\mathrm{T}}
    {U}{\Upsilon}
    %{V}{}
    {W}{\Omega}
    {X}{\Xi}
    {Y}{\Psi}
    {Z}{\mathrm{Z}}
   }
 }
\NewDocumentCommand{\lcgreek}{m}
 {
  \str_case:nn { #1 }
   {
    {a}{\alpha}
    {b}{\beta}
    {c}{\varsigma}
    {d}{\delta}
    {e}{\varepsilon}
    {f}{\varphi}
    {g}{\gamma}
    {h}{\eta}
    {i}{\iota}
    {j}{\vartheta}
    {k}{\kappa}
    {l}{\lambda}
    {m}{\mu}
    {n}{\nu}
    {o}{o}
    {p}{\pi}
    {q}{\chi}
    {r}{\rho}
    {s}{\sigma}
    {t}{\tau}
    {u}{\upsilon}
    %{v}{}
    {w}{\omega}
    {x}{\xi}
    {y}{\psi}
    {z}{\zeta}
   }
 }
\ExplSyntaxOff

\newcommand{\alphabetname}[1]{\uppercase{\ucgreek{#1}}}
